\subsection{Metodika a~ciele testovania}
\label{subsec:metodika}

Cieľom testovania bolo overiť, či implementovaná
aplikácia spoľahlivo rozpoznáva navrhnuté gestá
za podmienok, ktoré sa približujú reálnemu použitiu
počas zápasu. Konkrétne sme si stanovili tri
testovacie ciele:

\begin{enumerate}
  \item Overiť, či hraničné hodnoty stanovené v~kapitole
        \ref{chapter:gesta} zabezpečujú spoľahlivé
        rozpoznanie každého gesta pri jeho izolovanom
        vykonaní.
  \item Overiť, či kompletné kompozitné sekvencie
        (napríklad bodovanie s~viacerými flickami,
        napomenutie s~následným bodovaním súpera) prebiehajú
        bez chýb a~vedú k~očakávanému stavu aplikácie.
  \item Overiť odolnosť aplikácie voči falošne pozitívnym
        detekciám pri bežnej aktivite ruky, ktorá počas
        zápasu sprevádza prácu rozhodcu.
\end{enumerate}

\subsubsection{Postup testovania}

Testovanie sme rozdelili do troch fáz, ktoré
postupne zvyšujú zložitosť a~približujú sa reálnym
podmienkam použitia. Každá fáza zodpovedá jednému
z~vyššie uvedených cieľov:

\begin{itemize}
  \item \textbf{Fáza 1 -- jednotlivé gestá} -- každé gesto
        bolo testované izolovane v~stoji, v~kontrolovanom
        prostredí. Pre každé gesto sme vykonali 20 pokusov.
  \item \textbf{Fáza 2 -- kompozitné sekvencie} -- testovanie
        kompletných scenárov, ktoré simulujú reálne použitie
        aplikácie počas zápasu. Pre každý scenár sme
        vykonali 5 pokusov.
  \item \textbf{Fáza 3 -- falošne pozitívne detekcie} --
        10-minútový test pri bežnej aktivite ruky bez
        úmyslu vyvolať akékoľvek gesto.
\end{itemize}

\subsubsection{Podmienky testovania}

Všetky testy vykonával jeden používateľ -- autor
diplomovej práce. Hodinky boli umiestnené na ľavom
zápästí v~štandardnej polohe, identicky ako pri
zbere dát v~sekcii~\ref{sec:zber-dát}. Testovanie
prebiehalo v~kontrolovanom prostredí (uzavretá
miestnosť, bez sprievodných osôb), aby výsledky
neovplyvňovali externé faktory.

Aplikácia bola spustená v~debug móde
(viď podsekciu~\ref{app:debug-final}), ktorý
zobrazuje grafy senzorových dát, aktuálny režim
a~rozpoznané gestá v~reálnom čase. Tento mód
umožňuje okamžitú vizuálnu kontrolu, či bolo gesto
správne rozpoznané, prípadne identifikáciu príčiny
zlyhania detekcie.

Medzi jednotlivými pokusmi bola ruka vrátená do
neutrálnej polohy pri tele a~aplikácia bola
v~režime \texttt{WAITING}. Pred každým pokusom
sme overili, že aplikácia je v~očakávanom
počiatočnom stave.

\subsubsection{Hodnotené metriky}

Pre vyhodnotenie testovania sme použili nasledujúce
metriky:

\begin{itemize}
  \item \textbf{Úspešnosť rozpoznania} -- podiel správne
        rozpoznaných gest z~celkového počtu vykonaných
        pokusov, vyjadrený v~percentách. Pokus sa
        považuje za úspešný iba ak aplikácia rozpoznala
        práve to gesto, ktoré sme vykonali, a~nedošlo
        k~chybnej detekcii iného gesta.

  \item \textbf{Úspešnosť kompozitných sekvencií} --
        podiel sekvencií, pri ktorých výsledný stav
        aplikácie (skóre, aktívny režim, zaznamenané
        udalosti) zodpovedal očakávaniu. Sekvencia
        sa hodnotí ako úspešná iba ak prebehli všetky
        jej kroky správne -- ak zlyhalo čo i len jedno
        gesto, celá sekvencia sa považuje za neúspešnú.

  \item \textbf{Falošne pozitívne detekcie} -- počet
        rozpoznaných gest v~situáciách, keď žiadne
        gesto nebolo úmyselne vykonané. Pre každú
        falošnú detekciu sme zaznamenali aj typ
        nesprávne rozpoznaného gesta a~činnosť, pri
        ktorej k~detekcii došlo.
\end{itemize}

\subsubsection{Spôsob zaznamenávania výsledkov}

Pre každý vykonaný pokus sme zaznamenali, či bolo
gesto správne rozpoznané, a~v~prípade chyby aj
stručný popis príčiny (napríklad nesprávne rozpoznanie
ako iné gesto, žiadna detekcia, predčasná detekcia).
Tieto pozorovania sme následne využili pri diskusii
výsledkov a~identifikácii limitácií systému.